\section{Parameter values from 1-dimensional models}
\label{sec:5.1}

\begin{figure}
\centering
\staticfigure[.8\textwidth]{figures/chapter5/fig5.1a-f.png}{Posterior PDFs for each parameter in the single-parameter models (not marginalised over $H_0$). Each data set is denoted by a different colour, with the $1\sigma$ credible region in the lightest shade and $3\sigma$ the darkest. Regions in black are outside $3\sigma$. Larger versions of the $1\sigma$ credible regions are contained in \cref{app:large-figures}. (b) shows the Hubble parameter $h/h_{72}$, (c) the matter density $\Omega_m$ ; (d) the constant equation of state parameter $w_0$; (e) the varying equation of state parameter $w_1$; and, (f) the mixing factor $\gamma$.}
\end{figure}

The single-parameter models represent the simplest modification to the concordance $\Lambda$CDM model, containing only matter (both baryonic and dark) and a cosmological constant. Each posterior represents the probability that a value from the continuum set by the prior is the value which created the data. Thus, the mode is the value with the highest likelihood of being the value in the observable universe. The error is taken to be the width of the 68\% credible region by convention. The results are summarised in \cref{tab:5.1}. Plotting the credible regions themselves (\cref{fig:fig5.1a-f}) allows a comparison of the 1$\sigma$, 2$\sigma$, 3$\sigma$ errors for each data set.
   
The first step is to determine the value of $H_0$ used by the different SNe observing groups. This serves two purposes: it acts as a plausibility check on the smaller-error data (since these sets should agree (i.e. their 1$\sigma$ intervals should overlap) with the $H_0$ value of the sets on which they are based) and it justifies the use of the Hubble parameter as a free parameter for all models. If the $H_0$ value were the same for all data sets, I would have been able to set it as a constant, thus reducing the dimensionality of the model parameter spaces and shortening the length of the Markov chains. We find in \cref{fig:fig5.1a-f}(b) and \cref{tab:5.1} that only between the Gold and Constitution sets is there overlap in the 3$\sigma$ regions and that the credible regions of each reduced-error set are entirely enclosed by the same regions in the corresponding original set. Therefore we can conclude that a different value for the Hubble constant was applied to each normal-error set and that the reduced-error sets are behaving plausibly. It is vital to stress that I have not discovered a value for the Hubble constant: these values were applied arbitrarily by the different SNe search teams and then marginalised over as a nuisance parameter. This is in contrast to the other credible regions, which do suggest a value for their parameters.

The single-parameter models indicate that only minor modifications to the concordance model produce a better fit to the observed data. All three posteriors from the observational data give parameter values with consistent agreement, i.e. the 1$\sigma$ confidence regions overlap. 

The cosmological matter density $\Omega_m$ (\cref{fig:fig5.1a-f}(c)) has a best-fit value of $0.29$ across the Constitution and Union sets and their reduced-error sets and this agreement is confirmed by the overlap of the $1\sigma$ regions. The Gold set $1\sigma$ region intersects the $1\sigma$ regions of the Constitution and Union sets at $\Omega_m = 0.31$, however one must look at the 2$\sigma$ confidence regions before an agreed value is found between these and the Gold 0.2x set. The Gold and Gold 0.2x sets have a best-fit parameter value of $\Omega_m$ = 0.34, which is ruled out at 3$\sigma$ for the Constitution and Union sets and 4$\sigma$ for the corresponding 0.2x error sets. This demonstrates that the smaller-error sets amplify the calibration differences between the Gold set and the later Constitution and Union sets.

Similarly, the constant parameter $w_0$ (\cref{fig:fig5.1a-f}(c)) in the dark energy equation of state is well-agreed upon by all three observational data sets, which display overlap in their $1\sigma$ credible regions at $w_0 = -1$, although one must go to $2\sigma$ for overlap between the $1\sigma$ Gold 0.2x region and either the Constitution or Union sets. Once again the Gold 0.2x set does not overlap with the other two reduced-error sets. The agreement between the Constitution and Union samples is not as clear as for the matter density, as their best-fit values no longer agree, nor is there overlap at $1\sigma$ between their reduced-error sets. Nevertheless, the considerable overlap between the normal-error sets suggests that a cosmological constant $w_0 = -1$ can produce the observed SNe data.
   
There is no clear value for a time-varying dark energy equation of state which would better explain the distance moduli from the SNe data. This is once again due to the lack of consistency between the Gold set and the Union and Constitution sets. The latter two sets and their reduced-error sets all have a best-fit value around $w_1 = -0.25$ \cref{tab:5.1} whereas the Gold and Gold 0.2x sets share a best-fit value of $w_1 = 0.73$. The Gold best-fit value is ruled out at $3\sigma$ by the other two sets, whose best-fit value is slightly more plausible, lying in the 2$\sigma$ region for the Gold set. The 1$\sigma$ regions for all three observed sets intersect at $w_1 = 0.2$ in \cref{fig:fig5.1a-f}(e) whereas the concordance value of $w_1 = 0$ is in the 2$\sigma$ region for the Gold set. This suggests the inclusion of a small, positive variation in the equation of state. Given the difference in best-fit values between the Gold set and either the Constitution and Union sets, any modification would need to be tested against a sample of SNe at higher redshift (when the time-varying parameter has greater effect) to ascertain which of $w_1 \sim -0.25$, $w_1 \sim 0.2$ or $w_1 \sim 0.73$ is the most suitable value.
   
The interaction parameter $\gamma$ in \cref{fig:fig5.1a-f}(f) presents the same dilemma: although there is overlap of the 1$\sigma$ credible regions for the Gold and Union sets (the Constitution and Gold sets intersect at their 1 - 2$\sigma$ boundaries), the best-fit values lie in the 2$\sigma$ region of the other set. The considerable overlap between the Constitution and Union sets and their respective reduced-error sets is again present, all of which prefer (Table 5.1) a small positive interaction parameter to none at all. In contrast the Gold and Gold 0.2x sets prefer a negative interaction parameter approximately one order of magnitude greater than the Union and Constitution sets’ best-fit value. Thus the observational data propose a contradicting energy transfer: $\gamma > 0$ represent the transfer of energy from dark energy into matter, whereas $\gamma < 0$ represents the opposite. There is no clear value which satisfies all observational sets.

The projected sets from the \snap{} survey were artificially generated, so they do not propose parameter values for the observable Universe. As with the $H_0$ credible regions, we are interested solely in the reproduction of the values used to generate the distance moduli and the width of the region compared to the other data sets. From \cref{tab:5.1} we see that both the normal and reduced-error sets display the same best-fit values: moreover, these correspond to the concordance model values which were used to generate the data. This demonstrates that the Markov chains were run for a sufficient time to achieve convergence with the true posterior (i.e. stable equilibrium). Although the Snap 0.2x set display 1$\sigma$ regions which are far narrower than the other data sets, the Snap errors are of the same order of magnitude as those of the reduced-error sets. This suggests that a recalibration of existing data which reduces fivefold the uncertainty in distance moduli may produce equally tight constraints on the cosmological parameters as the pessimistic view of results from a new survey.

Comparison to the original papers in \cref{tab:5.2} is reassuring: the posteriors do not contradict the existing results for the Union and Constitution sets. Unfortunately, the tension in single-parameter models between the Gold set and these two sets is not reproduced in the original papers. This is because their value for the matter density was fixed at $\Omega_m = 0.28 \pm 0.04$, rather than the 0.3 which I used, so their constraints on the dark energy equation of state cannot be readily compared (\cite{Riess2007}). Their decision to use $\Omega_m$ as a Gaussian prior with maximum at 0.28 and width 0.04 came from large-scale structure measurements to obtain $\Omega_m$ h and Cepheids (another type of standard candle) to find a value for $H_0$ \cite{Riess2007}, so the validity of their results are sensitive to changes in the accepted values for $H_0$ . Therefore it is difficult to obtain a meaningful comparison between my single-parameter results and those of Riess et al. The 0.2x sets produce narrower confidence intervals due to the scaling of the error bars and the same best-fit values because they share the same data points as the observed data. Although this situation is less than ideal because the errors which contribute to the $\mu(z)$ uncertainties also produce scatter in the data points, there was no information available to guide any ”correction” to this scatter, so alterations would have generated meaningless data. The reproduction of best-fit values acts as a ”sanity check” to confirm convergence of the Markov chains. Quantitative attempts to compare how the parameter uncertainties scale with those in the distance moduli are limited by the resolution of the binning. Nevertheless, for both the normal- and reduced-error sets, my best-fit results do fit within and the 1$\sigma$ limits do overlap with (where comparison is appropriate) the 1$\sigma$ limits quoted in the papers.
